\subsection{Inter-Judge Reliability}
\label{app:interjudge-reliability}

The choice--score consistency reported above is an internal-consistency
measure: it asks whether a judge's selected winner is also the response to
which that judge assigned the higher mean rubric score. It does not establish
that different judges select the same winner. We therefore separately assess
inter-judge reliability on output comparisons evaluated by at least two
eligible judges.

The leave-family-out design creates a partially crossed panel because a judge
is ineligible whenever a comparison contains an output from its own model
family. We consequently report both raw agreement and Krippendorff's nominal
alpha, which accommodates missing judge--comparison cells and adjusts for
chance agreement. Before calculation, we map each choice to a canonical
output-pair order so that randomized presentation as Option~1 or Option~2
cannot affect agreement. We retain one evaluation per judge and comparison,
preventing comparisons with repeated evaluations from receiving additional
weight.

Across 6,265 output comparisons evaluated by at least two eligible judges,
the panel generates 11,257 judge-pair decisions. Judges select the same winner
in 71.0\% of these decisions, with Krippendorff's nominal
$\alpha=0.414$. Reliability is higher in automation
(74.5\%; $\alpha=0.448$) than in augmentation
(67.8\%; $\alpha=0.354$). Thirty-five comparisons with fewer than two valid
ratings are excluded from the reliability calculation.

\begin{figure}[H]
    \centering
    \includegraphics[width=\linewidth]
    {paper_figures/figureA_interjudge_reliability.pdf}
    \caption{
    \textbf{Inter-judge reliability across ten independent runs.}
    Panel A reports raw agreement and Krippendorff's nominal alpha overall and
    separately for augmentation and automation. Raw agreement is the share of
    eligible judge-pair decisions selecting the same canonical winner; alpha
    adjusts for chance agreement and accommodates missing ratings induced by
    leave-family-out evaluation. Panel B reports alpha by task and usage mode.
    Only comparisons evaluated by at least two eligible judges are included.
    }
    \label{fig:interjudge-reliability}
\end{figure}

Task--mode reliability varies substantively. Krippendorff's alpha ranges from
0.231 for operations-research augmentation to 0.615 for market-trends
automation. Automation produces higher agreement than augmentation overall
and in every task shown in Figure~\ref{fig:interjudge-reliability}. These
results indicate a meaningful but moderate common signal across judges rather
than near-consensus. The lower agreement in augmentation also suggests that
evaluating the downstream effects of assistance texts is more
judge-dependent than evaluating direct task outputs. We therefore interpret
aggregated rankings as panel estimates that reduce idiosyncratic judge effects,
not as ground-truth labels. This finding further motivates the human-expert
validation proposed in the limitations section.
